% Generated draft: requires mathematical and source review.
\documentclass[11pt]{article}
\usepackage[margin=25mm]{geometry}
\usepackage{fontspec}
\setmainfont{Latin Modern Roman}
\usepackage{amsmath,amssymb,mathtools,unicode-math}
\setmathfont{Latin Modern Math}
\usepackage{enumitem,needspace,xurl,hyperref}
\setlength{\parindent}{0pt}
\setlength{\parskip}{5pt}
\setlength{\emergencystretch}{3em}
\newcommand{\N}{\mathbb{N}}
\newcommand{\Z}{\mathbb{Z}}
\newcommand{\Q}{\mathbb{Q}}
\newcommand{\R}{\mathbb{R}}
\newcommand{\C}{\mathbb{C}}
\newcommand{\F}{\mathbb{F}}
\newcommand{\A}{\mathbb{A}}
\newcommand{\PP}{\mathbb{P}}
\newcommand{\E}{\mathbb{E}}
\newcommand{\eps}{\varepsilon}
\newcommand{\dd}{\,\mathrm{d}}
\begin{document}
\section*{UnsolvedMath ID: 30005317}
\textbf{Source problem code:} catalogue entry
\subsection{Definitions and mathematical statement}
A domino is a pair of unit lattice squares sharing an edge. A domino tiling of a lattice region is a partition of its unit squares into dominoes. All configuration and tiling counts below are unweighted; equivalently, every allowed configuration has weight \(1\).

The supplied record concerns the twenty-vertex model on a triangular lattice. In the usual local interpretation, each interior vertex has six incident edges, and an allowed orientation has three edges directed inward and three directed outward, giving twenty local possibilities. The record describes its boundary condition DWBC3 by imposing six-vertex domain-wall boundary conditions on the horizontal and vertical boundary edges and directing all diagonal boundary edges up/left. It does not specify the individual horizontal and vertical boundary arrows, the precise finite lattice domain, or the convention for counting its rows and columns. Consequently, the following sets refer to the source's intended domains and boundary conventions, which require clarification before this becomes a fully explicit standalone statement.

Let \(\mathcal C_n\) denote the set of twenty-vertex configurations with DWBC3 on the source's \(n\times n\) domain. Let \(P_n\) denote Pachter's triangle of size \(2n\), described in the record as obtained by halving a \(2n\times 2n\) square along a diagonal staircase whose steps have length \(2\). Write \(\mathcal T(P_n)\) for its domino tilings. This description does not fix the staircase placement or the retained half sufficiently to give a coordinate definition here.

For a nonnegative integer \(k\), let \(\mathcal C_{n,k}\) denote the configurations of the source's extended twenty-vertex model on an \((n+k)\times n\) rectangle, with the extended DWBC3 boundary condition and horizontal outgoing edge orientations mentioned in the record. The record does not identify all affected boundary edges. Let \(P_{n,k}\) denote the corresponding enhanced Pachter triangle, obtained by raising by \(k\) the ``roof'' constraining its nonintersecting lattice-path formulation. Neither that path formulation nor the roof coordinates and the reconstruction of the enhanced domino region are supplied.

Let \(\mathcal I\) denote the intended admissible positive integer sizes \(n\); the record does not explicitly give their range. Subject to resolving the preceding definitions, the two source conjectures ask whether
\[
  \lvert\mathcal C_n\rvert=\lvert\mathcal T(P_n)\rvert
  \qquad\text{for every }n\in\mathcal I
  \tag{Conjecture 3}
\]
and
\[
  \lvert\mathcal C_{n,k}\rvert=\lvert\mathcal T(P_{n,k})\rvert
  \qquad\text{for every }n\in\mathcal I
  \text{ and every }k\in\mathbb Z_{\geq 0}
  \tag{Conjecture 4}
\]
hold. The integer interpretation of \(k\) reflects its role as an increment in the number of grid columns. The extension is presented as encompassing the square case at \(k=0\), but the identities \(\mathcal C_{n,0}=\mathcal C_n\) and \(P_{n,0}=P_n\) must be checked against the missing domain definitions. In particular, the universal quantifier over \(k\) is retained from the record; it is not replaced by a bounded range.

The source calls \(k=n-1\) the maximal raised-roof enhancement. It reports that the corresponding configuration numbers, denoted there by \(Z_n^{20Ve}\), were obtained in terms of the numbers \(Z_n^{6VU}\) for a six-vertex model with U-turn boundary conditions. The passage also refers to coefficients of a \(u,v\) expansion in its Theorem 1, matrix entries of \(I+M\), and diagonal minors counting nonintersecting paths with horizontal, diagonal, and vertical steps. Those paths are said to be in bijection with domino tilings of a conic fundamental domain of a holey Aztec square under quarter-turn symmetry. The theorem, matrix, domains, and explicit relation between the two sequences are absent from the supplied record, so this passage does not supply an additional usable formula here. The meaning of ``maximal'' at \(k=n-1\) needs reconciliation with the stated range of all \(k\geq0\).

For the uniformly weighted square-domain partition function \(Z_n^{20V3}=\lvert\mathcal C_n\rvert\), the record lists
\[
\begin{gathered}
1,\ 3,\ 29,\ 901,\ 89893,\ 28793575,\ 29607089625,\\
97725875584681,\ 1035449388414303593,\ \ldots .
\end{gathered}
\]
The starting index is not explicitly supplied. The record also says that this partition function has no direct connection to the six-vertex partition function \(Z_n^{6V}\); it does not define the intended meaning of ``direct connection.'' These are supplied contextual assertions, not proofs of either conjecture.

The supplied dated literature assessment labels the problem partially solved and reports proofs of the twenty-vertex/domino correspondence for the maximal domain and the \(k=2,3\) truncations, while saying that the square-grid and all-\(k\) assertions are not fully proved. The meaning of those truncations and their precise relationship to \(\mathcal C_{n,k}\) and \(P_{n,k}\) are not specified. This assessment is retained as source curation and is not independently verified.
\subsection{Short English statement}
The two conjectures compare unweighted counts of twenty-vertex configurations with DWBC3 and domino tilings of Pachter's triangle, first on a square domain and then on an extended rectangle with a roof raised by every integer \(k\geq0\). Exact boundary arrows, region coordinates, and the roof construction remain unspecified. The supplied partial-solution assessment is dated curation, not independent verification. This statement-only draft requires human mathematical review.
\subsection{Sources}
\begin{itemize}
\item Local UnsolvedMath registry, permanent ID \(30005317\), title ``Twenty-Vertex Models and Pachter-Triangle Tilings,'' source problem code \texttt{catalogue entry}, in the Oberwolfach Reports Open Problems collection. The supplied record includes Conjectures 3 and 4 and contextual material. No proposer is named; the proposed year is recorded as 2022.
\item Supplied report citation: \emph{Enumerative Combinatorics (2023)}, Oberwolfach Reports, DOI \texttt{10.4171/owr/2022/57}, pp.\ 3215--3305. Supplied source URL and literature-source link labelled ``original OWR report'': \url{https://doi.org/10.4171/owr/2022/57}. The record describes Oberwolfach Reports as documenting mathematical workshops and their open-problem discussions.
\item Supplied literature-source link labelled ``Di Francesco primary article'': \url{https://www.combinatorics.org/ojs/index.php/eljc/article/download/v28i4p38/pdf}. The record supplies no article title, full author list, or precise theorem locations. This label does not establish authorship of the conjectures.
\item Literature assessment checked at \texttt{2026-08-22}: supplied dated curation, not independently verified. It labels the problem partially solved, reports a proved correspondence for the maximal domain and the \(k=2,3\) truncations, and says that the square-grid and all-\(k\) assertions are not fully proved. The background expressly calls this a dated triage derived from the supplied OWR curation and requests independent verification of time-sensitive status and cited claims. Neither linked source has been visited for this draft.
\item The source excerpt uses reference labels [10], [13], and [14], attributing the rectangle extension to [10], the triangle to Pachter [13], and the U-turn six-vertex relation to [14]. Their bibliographic entries are not supplied. Its reference to Theorem 1 is likewise unresolved.
\item The registry's statement-status label is \texttt{corrected\_verified}; its accompanying explanation limits the repairs to formula delimiters, notation, grammar, and typography and says that mathematical content and scope were preserved. This label is not independent mathematical verification. The record was updated at \texttt{2026-09-16T12:06:53Z}.
\end{itemize}
\end{document}
